\section{Distilling the Context Generator}
\label{sec:distill}

The teacher produces high-quality contexts but requires an API call per chunk
and sends document text to a third party. We therefore distill the teacher into
\texttt{Qwen3-0.6B} \citep{qwen3} so that the ingest stage runs locally.

\paragraph{Training.}
The student is fine-tuned with QLoRA \citep{dettmers2023qlora}: the base model
is loaded in 4-bit precision, and low-rank adapters \citep{hu2022lora} with
rank 16 and scaling factor 16 are attached to all linear layers, which
corresponds to 1.665\% of the parameters (10.1M of 606.1M). Training uses two
epochs with an effective batch size of eight, a learning rate of
$2\times10^{-4}$, and a maximum sequence length of 16{,}384 tokens; the longest
training sample is 9.7k tokens, so no sample is truncated. Loss is computed on
the assistant response only, with the document and the instruction masked out.
The training and inference prompts are byte-identical
chat templates with thinking disabled, which is enforced by a hard check before
training.

\paragraph{Quality.}
Table~\ref{tab:distill} compares the base model and the fine-tuned student on
sixty held-out chunks, using word-level F1 against the teacher context as the
metric. The base model is unusable: it produces empty or unrelated output, with
a mean F1 of 0.189 and 1.7\% empty responses. After fine-tuning, the mean F1 is
0.537, the tenth percentile rises from 0.064 to 0.319, and no response is
empty. Neither model emits numbers that are absent from the source document, so
the improvement does not come from hallucinated detail. Generation is slightly
slower per sample (7.82~s against 5.55~s) because the fine-tuned model produces
longer, more specific contexts instead of short generic ones.

\begin{table}[t]
\centering
\small
\caption{Distillation quality: word-level F1 of generated contexts against the teacher, on 60 held-out chunks (base Qwen3-0.6B vs.\ the same model after LoRA SFT).}
\label{tab:distill}
\begin{tabular}{lrr}
\toprule
Metric & Baseline & After SFT \\
\midrule
word-F1 mean & 0.189 & \textbf{0.537} \\
word-F1 median & 0.176 & 0.534 \\
word-F1 p10 & 0.064 & 0.319 \\
Empty responses & 1.7\% & \textbf{0\%} \\
Hallucinated numbers & 0 & 0 \\
Generation time / sample (s) & 5.55 & 7.82 \\
\bottomrule
\end{tabular}
\end{table}


\paragraph{Corpus-wide ingest.}
We regenerate the context for all 3{,}134 chunks with the student on the target
device (Apple~M2, MLX in bfloat16). All chunks of a window share the same
prefix, so its key/value state is computed once and reused, which reduces
prefill from an estimated 7.08~h to 2.31~h, a factor of 3.1. Total wall-clock
time is 4.49~h, split into 2.31~h of prefill and 2.19~h of decode; the median
chunk takes 4.20~s. Eleven contexts (0.35\%) reach the 160-token output cap.
The teacher required 1.85~h and USD~0.82, so the on-device path is 2.4 times
slower, costs nothing per run, and never transmits document text.
Table~\ref{tab:ingest} reports the comparison.
